(1) 两质点速度的一阶展开给 \(T=m\ell^2(2\dot\theta_1^2+2\dot\theta_1\dot\theta_2+\dot\theta_2^2)/2\)；相对最低位置的势为 \(V=mg\ell(2\theta_1^2+\theta_2^2)/2\)。故质量矩阵 \(M=m\ell^2\begin{pmatrix}2&1\\1&1\end{pmatrix}\)，刚度矩阵 \(K=mg\ell\,\mathrm{diag}(2,1)\)。交叉动能源于下方质点同时接受两杆的水平速度。
(2) 令 \(\theta=u e^{i\omega t}\)，则 \(\det(K-\omega^2M)=0\)，记 \(\lambda=\omega^2\ell/g\) 得 \(\lambda^2-4\lambda+2=0\)。因此 \(\omega_\pm^2=(g/\ell)(2\pm\sqrt2)\)。从第一行 \(2(1-\lambda)u_1-\lambda u_2=0\) 得低频模 \(u_-=(1,\sqrt2)\)，高频模 \(u_+=(1,-\sqrt2)\)，二者按 \(M\) 内积正交。
(3) 把初始向量 \((0,\theta_0)\) 展为 \(A u_-+B u_+\)，由 \(A+B=0\)、\(\sqrt2(A-B)=\theta_0\) 得 \(A=\theta_0/(2\sqrt2)\)、\(B=-A\)。初速度全零，因此 \(\theta_1(t)=A[\cos(\omega_-t)-\cos(\omega_+t)]\)，\(\theta_2(t)=\theta_0[\cos(\omega_-t)+\cos(\omega_+t)]/2\)，代入可直接恢复两个初始角和零速度。
(4) 初始动能为零、势能为 \(E_0=mg\ell\theta_0^2/2\)。将模态解代回 \(T+V\)，利用 \(u_-^TMu_+=0\) 及 \(Ku_j=\omega_j^2Mu_j\)，两个模能量分别守恒且总和恰为 \(E_0\)。结果只在 \(|\theta_{1,2}|\ll1\) 下准确；两摆角独立并不代表动能可拆开，因为第二质点受第一个铰的平动耦合。
